\chapter{Discussion}
\label{ch:discussion}

This chapter interprets the descriptive and inferential patterns
reported in Chapter~\ref{ch:evaluation} against the two research
questions of this thesis. The numerical results themselves are not
restated here; the reader is referred to
Section~\ref{sec:eval-summary} for the consolidated recap.

\section{RQ1: User Experience Across Modes}
\label{sec:disc-rq1}

On the self-report instruments, the structured Form mode is rated
the most usable, the least demanding, and the most preferred.
Form receives the largest share of first-place preference votes
($53.6\,\%$), the highest aggregate SUS-item ratings, and the lowest
means on every R-TLX demand subscale. The Friedman omnibus tests on
the SUS and R-TLX measures do not reach the conventional significance
threshold; the significant sub-effects on the matched-pairs subset
are confined to Form versus the LLM-driven modes on Temporal Demand,
Frustration, and re-use intent.

A pattern that recurs across every aggregate user-experience measure
is that \emph{Ask and Chat are descriptively and inferentially
indistinguishable} from each other: pairwise tests between the two
LLM-driven modes are non-significant on every SUS aggregate, every
R-TLX subscale, every pooled AI-support item, and the rank
distribution. The user-experience contrast in this study is therefore
best read as a binary one between the structured questionnaire (Form)
and the LLM-driven interactions (Ask, Chat together), rather than as
the three-way ordering implied by the original ``more LLM = better''
framing.

The Perceived AI Support items do show one monotonic pattern with
LLM involvement: legal-basis confidence (item~$5_{ai}$) rises from
Form to Chat. The same items show a counter-pattern on re-use intent
(item~$6_{ai}$), where Form is rated highest. The free-text comments
characterise Form as fast and ergonomic, Ask as transparent but
cumbersome, and Chat as comfortable but opaque, which is consistent
with the quantitative split.

\section{RQ2: Document Quality and the Confidence Gap}
\label{sec:disc-rq2}

The document-quality side of the study yields a more nuanced ordering
than the self-report side. Reference-based NLP metrics place Form
first on lexical recall (BERTScore Recall, ROUGE-2) and the n-gram-
based metrics, Chat first on semantic similarity (SBERT) and BLEU,
and Ask consistently last on every metric. The chapter-level
breakdown attributes Chat's BLEU and METEOR advantage primarily to
the Retention Periods and the Technical and Organisational Measures
chapters.

The LLM-as-Judge inverts the BERTScore Recall ordering. Form leads
on Completeness ($4.84$) but lags on Scenario Faithfulness ($1.84$)
and Hallucination (inverse, $1.65$); Chat leads on Faithfulness and
(inverse) Hallucination but trails on Completeness; Overall scores
are clustered tightly between $2.41$ and $2.62$. The proximate source
of this inversion is the joint distribution noted in
Section~\ref{sec:eval-summary}: $83.8\,\%$ of Form documents combine
maximum Completeness with moderate-to-significant hallucination, an
artefact of the checkbox interface that produces structurally
complete but scenario-detached documents. This mixed picture sits
uneasily with the implicit ``more LLM equals better quality''
assumption that motivates the architecturally elaborated pipelines
of recent legal-document generation work \cite{nigamStructuredLegalDocument2025,linLegalDocumentsDrafting2024,liCaseGenBenchmarkMultiStage2025}.

Reading RQ1 and RQ2 together surfaces an apparent gap between
\emph{perceived} AI support and \emph{objective} document quality.
Users rate Form as the most usable and prefer it most, while the
mode that most often matches the reference at the semantic level
(Chat) is rated as opaque and ranked second. The mode that scores
highest on AI-mediated legal-basis confidence (Chat) is also the
mode that produces the documents the judge rates as least complete.
This parallels Surya et al.'s evaluation of an interactive legal AI
assistant \cite{suryaAIPoweredInteractiveLegal2025}, where ease of
use was the headline finding — a pattern that here translates into
the structured Form being preferred despite its weaker faithfulness
profile. The implications of this gap for the design of LLM-supported
ROPA tooling, and the extent to which the present novice sample
generalises, are addressed in the remainder of this chapter.
